\section*{Chapitre \ref{ch:g10:universal-gravitation} --- Gravitation universelle et poids}

\begin{solution}{exo:g10:universal-gravitation:1}
\emph{1.} $F = \num{6.67e-11} \times \dfrac{60 \times 60}{0.80^2}
= \dfrac{\num{2.40e-7}}{0.64} \approx \qty{3.8e-7}{N}$.

\emph{2.} Le \omterm{def:g10:universal-gravitation:weight}{poids} du moustique est
$P = \num{2.5e-6} \times 9.81 \approx \qty{2.5e-5}{N}$ --- environ
$65$~fois \emph{plus grand} que l'attraction mutuelle des danseurs.
La gravité entre personnes est plus faible que le \omterm{def:g10:universal-gravitation:weight}{poids} d'un
moustique.
\end{solution}

\begin{solution}{exo:g10:universal-gravitation:2}
\emph{1.} $P = mg = 55 \times 9.81 \approx \qty{5.4e2}{N}$.

\emph{2.} Sa masse est inchangée~: \qty{55}{kg} (même quantité de
matière). Son \omterm{def:g10:universal-gravitation:weight}{poids} devient
$P = 55 \times 1.62 \approx \qty{89}{N}$ --- six fois moins.
\end{solution}

\begin{solution}{exo:g10:universal-gravitation:3}
\emph{1.} $g_{\text{Mars}} = \dfrac{G M}{R^2}
= \dfrac{\num{6.67e-11} \times \num{6.42e23}}{(\num{3.39e6})^2}
= \dfrac{\num{4.28e13}}{\num{1.15e13}} \approx \qty{3.73}{N/kg}$.

\emph{2.} Sur Mars~: $P = 900 \times 3.73 \approx \qty{3.4e3}{N}$. Sur
Terre~: $P = 900 \times 9.81 \approx \qty{8.8e3}{N}$ --- le rover
pèse environ $2{,}6$~fois moins sur Mars, avec la même masse.
\end{solution}

\begin{solution}{exo:g10:universal-gravitation:4}
\[
F = \num{6.67e-11} \times
\frac{\num{5.97e24} \times \num{7.35e22}}{(\num{3.84e8})^2}
= \frac{\num{2.93e37}}{\num{1.47e17}}
\approx \qty{2.0e20}{N}.
\]
La Lune exerce sur la Terre une force de la \emph{même} intensité,
\qty{2.0e20}{N}, en sens opposé~: la \omterm{def:g10:universal-gravitation:force}{gravitation} est mutuelle.
\end{solution}

\begin{solution}{exo:g10:universal-gravitation:5}
\emph{1.} Distance doublée~: $d^2$ est multiplié par $4$, donc
$F = F_0/4$.

\emph{2.} Distance divisée par deux~: $F = 4F_0$.

\emph{3.} Les deux masses doublées~: $m_1 m_2$ est multiplié par $4$,
donc $F = 4F_0$.

\emph{4.} Numérateur $\times 3$, dénominateur $\times 9$~:
$F = F_0/3$.
\end{solution}

\begin{solution}{exo:g10:universal-gravitation:6}
\emph{1.} $F = \num{6.67e-11} \times
\dfrac{(\num{5.0e8})^2}{100^2}
= \num{6.67e-11} \times \dfrac{\num{2.5e17}}{\num{1.0e4}}
\approx \qty{1.7e3}{N}$.

\emph{2.} $m = F/g = 1667/9.81 \approx \qty{1.7e2}{kg}$~: les deux
navires s'attirent avec le \omterm{def:g10:universal-gravitation:weight}{poids} d'une grosse moto.

\emph{3.} $a = F/m = \num{1.7e3}/\num{5.0e8} \approx
\qty{3.3e-6}{m/s^2}$. Même sans opposition, les pétroliers mettraient
environ cinq jours à atteindre une allure de marche~; en pratique, la
résistance de l'eau et les amarres noient complètement l'attraction.
\end{solution}

\begin{solution}{exo:g10:universal-gravitation:7}
\emph{1.} La balance à plateaux indique \qty{3.0}{kg} sur les deux
mondes~: le sac et les masses étalons voient leurs
\omterm{def:g10:universal-gravitation:weight}{poids} divisés par le même facteur, donc la comparaison est
inchangée. Le dynamomètre lit le \omterm{def:g10:universal-gravitation:weight}{poids}~:
$3.0 \times 9.81 \approx \qty{29}{N}$ sur Terre,
$3.0 \times 1.62 \approx \qty{4.9}{N}$ sur la Lune.

\emph{2.} La balance~: elle mesure la masse, ce que le client achète.
Un dynamomètre calibré en «~\omterm{def:g10:orders-of-magnitude:unit}{kilogrammes}~» sur Terre lirait
$4.9/9.81 \approx \qty{0.50}{kg}$ pour le même sac sur la Lune --- le
commerçant donnerait six sacs pour le prix d'un.
\end{solution}

\begin{solution}{exo:g10:universal-gravitation:8}
\emph{1.} Distance au centre~:
$d = R + h = \num{6.37e6} + \num{4.0e5} = \qty{6.77e6}{m}$. Alors
$g = \dfrac{\num{6.67e-11} \times \num{5.97e24}}{(\num{6.77e6})^2}
= \dfrac{\num{3.98e14}}{\num{4.58e13}} \approx \qty{8.69}{N/kg}$.

\emph{2.} $8.69/9.81 \approx 89\%$ de la valeur en surface.

\emph{3.} La gravité est presque à pleine force là-haut~: les
astronautes flottent parce que la station et tout ce qu'elle contient
\emph{tombent ensemble} autour de la Terre (chute libre), non parce
que la gravité a disparu.
\end{solution}

\begin{solution}{exo:g10:universal-gravitation:9}
De $g = GM/R^2$,
\[
M = \frac{g R^2}{G}
= \frac{8.87 \times (\num{6.05e6})^2}{\num{6.67e-11}}
= \frac{\num{3.25e14}}{\num{6.67e-11}}
\approx \qty{4.87e24}{kg},
\]
environ $0{,}8$~fois la masse de la Terre --- Vénus est presque la
jumelle de la Terre en taille et en masse.
\end{solution}

\begin{solution}{exo:g10:universal-gravitation:10}
\emph{1.} Soleil~:
$F = \dfrac{\num{6.67e-11} \times \num{1.99e30}}{(\num{1.496e11})^2}
\approx \qty{5.9e-3}{N}$. Lune~:
$F = \dfrac{\num{6.67e-11} \times \num{7.35e22}}{(\num{3.84e8})^2}
\approx \qty{3.3e-5}{N}$.

\emph{2.} Le Soleil tire environ $\num{5.9e-3}/\num{3.3e-5} \approx 180$
fois plus fort.

\emph{3.} Le diamètre de la Terre est une \emph{fraction} bien plus
grande de la distance Terre--Lune ($\approx 3\%$) que de la distance
Terre--Soleil ($\approx 0{,}01\%$), donc la traction de la Lune change
bien plus de la face proche à la face lointaine du globe. Les marées
se nourrissent de cette différence --- et là la Lune l'emporte.
\end{solution}

\begin{solution}{exo:g10:universal-gravitation:11}
\emph{1.} $g' = \dfrac{G (2M)}{(2R)^2} = \dfrac{2}{4}\,\dfrac{GM}{R^2}
= \dfrac{g}{2} \approx \qty{4.9}{N/kg}$~: le rayon doublé l'emporte
sur la masse doublée.

\emph{2.} Il faut $\dfrac{GM}{R'^2} = 2\,\dfrac{GM}{R^2}$, donc
$R'^2 = R^2/2$ et $R' = R/\sqrt{2} \approx \qty{4.5e6}{m}$ --- environ
\qty{4500}{km}.
\end{solution}

\begin{solution}{exo:g10:universal-gravitation:12}
\emph{1.} Tractions égales par \omterm{def:g10:orders-of-magnitude:unit}{kilogramme}~:
$\dfrac{G M_{\text{Earth}}}{x^2} = \dfrac{G M_{\text{Moon}}}{(D-x)^2}$.
En croisant et en divisant par $G$~:
$\left(\dfrac{x}{D-x}\right)^2 = \dfrac{M_{\text{Earth}}}{M_{\text{Moon}}}$.

\emph{2.} $\dfrac{x}{D-x}
= \sqrt{\dfrac{\num{5.97e24}}{\num{7.35e22}}} = \sqrt{81.2} \approx
9.0$, donc $x = 9.0\,(D - x)$, c'est-à-dire $x = \dfrac{9.0}{10.0}\,D
\approx 0.90 \times \num{3.84e8} \approx \qty{3.5e8}{m}$. Le point
d'équilibre se trouve à $90\%$ du trajet vers la Lune~: la traction de
la Terre domine presque tout le voyage.
\end{solution}

\begin{solution}{exo:g10:universal-gravitation:13}
\emph{1.} En développant~: $R^2 + 2Rh + h^2 = R^2 + x^2$~; comme $h$
est minuscule devant $R$, on abandonne $h^2$~: $2Rh \approx x^2$, donc
$h \approx \dfrac{x^2}{2R}
= \dfrac{(\num{8.0e3})^2}{2 \times \num{6.37e6}}
= \dfrac{\num{6.4e7}}{\num{1.27e7}} \approx \qty{5.0}{m}$.

\emph{2.} En une seconde le projectile tombe de \qty{4.9}{m} ---
presque exactement les \qty{5.0}{m} dont le sol s'abaisse sur
\qty{8.0}{km}. Donc à $v \approx \qty{8.0}{km/s}$ la chute ne rattrape
jamais le sol~: le projectile \omterm{def:g10:universal-gravitation:satellite}{orbite}.

\emph{3.} $T = \dfrac{2\pi R}{v}
= \dfrac{2\pi \times \num{6.37e6}}{\num{8.0e3}}
\approx \qty{5.0e3}{s} \approx \qty{83}{min}$ --- proche des environ
$\qty{90}{min}$ des vrais \omterm{def:g10:universal-gravitation:satellite}{satellites} en
\omterm{def:g10:universal-gravitation:satellite}{orbite basse}, qui volent un peu plus haut.
\end{solution}

\begin{solution}{exo:g10:universal-gravitation:14}
\emph{1.} $g = \dfrac{\num{6.67e-11} \times \num{2.8e30}}
{(\num{1.2e4})^2} = \dfrac{\num{1.87e20}}{\num{1.44e8}}
\approx \qty{1.3e12}{N/kg}.$

\emph{2.} $P = 70 \times \num{1.3e12} \approx \qty{9.1e13}{N}$, contre
$\qty{6.9e2}{N}$ sur Terre~: environ $\num{1.3e11}$~fois plus. Aucune
structure faite d'atomes ne survit debout là-bas.

\emph{3.} $P = \num{1.0e-6} \times \num{1.3e12} \approx
\qty{1.3e6}{N}$. Sur Terre c'est le \omterm{def:g10:universal-gravitation:weight}{poids} de
$m = \num{1.3e6}/9.81 \approx \qty{1.3e5}{kg}$ --- un grain de sable
pesant autant qu'une locomotive de cent tonnes.
\end{solution}

\begin{solution}{exo:g10:universal-gravitation:15}
\emph{1.} La distance est $60$~fois plus grande, donc la traction par
\omterm{def:g10:orders-of-magnitude:unit}{kilogramme} est $60^2 = 3600$~fois plus faible.

\emph{2.} La chute en une seconde se met à l'échelle de la même façon~:
$4.9/3600 \approx \qty{1.4e-3}{m}$ --- environ \qty{1.4}{mm}.

\emph{3.} $T = 27.3 \times \qty{86400}{s} = \qty{2.36e6}{s}$, donc
$v = \dfrac{2\pi \times \num{3.84e8}}{\num{2.36e6}} \approx
\qty{1.02e3}{m/s}$, et $x = \qty{1.02e3}{m}$ en une seconde. La chute
est $h \approx \dfrac{x^2}{2d}
= \dfrac{(\num{1.02e3})^2}{2 \times \num{3.84e8}}
\approx \qty{1.4e-3}{m}$.

\emph{4.} Les deux nombres s'accordent~: la Lune tombe vers la Terre
chaque seconde exactement autant que la gravité en inverse du carré,
calibrée sur une pomme qui tombe, le prédit. Newton conclut que la
force qui retient la Lune \emph{est} la force qui fait tomber la pomme
--- la gravité est universelle.
\end{solution}

\begin{solution}{pb:g10:universal-gravitation:1}
\textbf{1.} $F = \num{6.67e-11} \times
\dfrac{\num{7.3e6} \times 55}{100^2} \approx \qty{2.7e-6}{N}$.

\textbf{2.} Le grain de sable pèse
$\num{1.0e-6} \times 9.81 \approx \qty{9.8e-6}{N}$~: toute la tour
Eiffel tire le touriste environ $3{,}7$~fois \emph{moins} que le
\omterm{def:g10:universal-gravitation:weight}{poids} d'un grain de sable.

\textbf{3.} (a) Distance $\times 10$~: force divisée par $100$.
(b) Les deux masses $\times 1000$~: force multipliée par $10^6$.
(c) Masses $\times 1000$ et distance $\times 1000$~:
$10^6 / 10^6$ --- la force est inchangée.

\textbf{4.} $F = \dfrac{\num{6.67e-11} \times \num{5.97e24} \times
1.00}{(\num{6.37e6})^2} \approx \qty{9.81}{N}$~: le
\omterm{def:g10:universal-gravitation:weight}{poids} d'un \omterm{def:g10:orders-of-magnitude:unit}{kilogramme} --- c'est exactement
l'intensité du champ $g = \qty{9.81}{N/kg}$.

\textbf{5.} La gravité est faible par \omterm{def:g10:orders-of-magnitude:unit}{kilogramme} mais n'a
jamais d'autre effet qu'\emph{attirer}, donc les tractions de tous
les $\num{5.97e24}$ \omterm{def:g10:orders-of-magnitude:unit}{kilogrammes} de la Terre s'additionnent
dans le même sens --- et une masse astronomique bat une masse
voisine de vingt-deux puissances de dix.

\textbf{6.} $F = \num{6.67e-11} \times \dfrac{158 \times 0.73}{0.23^2}
= \num{6.67e-11} \times \dfrac{115.3}{0.0529}
\approx \qty{1.5e-7}{N}$.

\textbf{7.} La petite sphère pèse
$0.73 \times 9.81 \approx \qty{7.2}{N}$ --- environ $\num{5e7}$~fois
l'attraction à détecter. Aucune balance à plateaux ne résout une
partie sur cinquante millions~; un fil de torsion fin, qui tourne de
façon visible sous des couples à l'échelle du nanonewton, le peut.

\textbf{8.} $G = \dfrac{F d^2}{m_1 m_2}
= \dfrac{\num{1.47e-7} \times 0.0529}{115.3}
\approx \qty{6.74e-11}{N.m^2/kg^2}$ --- à environ $1\%$ de la valeur
moderne \num{6.67e-11}, depuis un hangar de bois en 1798.

\textbf{9.} De $g = G M / R^2$~:
\[
M = \frac{g R^2}{G}
= \frac{9.81 \times (\num{6.37e6})^2}{\num{6.67e-11}}
\approx \qty{5.97e24}{kg}.
\]
Six millions de milliards de milliards de
\omterm{def:g10:orders-of-magnitude:unit}{kilogrammes}~: la Terre, pesée.

\textbf{10.} $V = \frac43 \pi R^3 = \frac43 \pi
(\num{6.37e6})^3 \approx \qty{1.08e21}{m^3}$, donc
$\rho = \num{5.97e24}/\num{1.08e21} \approx \qty{5.5e3}{kg/m^3}$ ---
deux fois la masse volumique de la roche de surface. L'intérieur doit
cacher quelque chose de bien plus dense que le granite~: la Terre a
un noyau lourd (fer).

\textbf{11.} $g_{\text{Moon}} = \dfrac{\num{6.67e-11} \times
\num{7.35e22}}{(\num{1.74e6})^2} = \dfrac{\num{4.90e12}}
{\num{3.03e12}} \approx \qty{1.62}{N/kg}$.

\textbf{12.} Hauteur $\propto 1/g$~:
$h = 0.50 \times \dfrac{9.81}{1.62} \approx \qty{3.0}{m}$ --- un saut
debout par-dessus un panier de basket, au ralenti.

\textbf{13.} $g_{\text{Mercury}} = \dfrac{\num{6.67e-11} \times
\num{3.30e23}}{(\num{2.44e6})^2} = \dfrac{\num{2.20e13}}
{\num{5.95e12}} \approx \qty{3.70}{N/kg}$ --- presque exactement le
\qty{3.73}{N/kg} de Mars. Mercure a moins de masse, mais son rayon
plus petit (un $R^2$ plus petit en bas) compense~: c'est un monde bien
plus dense.

\textbf{14.} La moitié de l'intensité du champ exige
$(R + h)^2 = 2 R^2$, donc $R + h = R\sqrt{2}$ et
$h = (\sqrt{2} - 1) R \approx 0.414 \times \num{6.37e6}
\approx \qty{2.6e6}{m}$ --- environ \qty{2600}{km} en altitude.

\textbf{15.} $\dfrac{G M_{\text{Earth}}}{d^2} = 1.62$ donne
$d = \sqrt{\dfrac{\num{6.67e-11} \times \num{5.97e24}}{1.62}}
= \sqrt{\num{2.46e14}} \approx \qty{1.57e7}{m} \approx 2.5$ rayons
terrestres. Dès un rayon et demi au-dessus du sol, la Terre ne tire
pas plus fort que la surface de la Lune.

\textbf{16.} $v = \dfrac{2\pi d}{T}
= \dfrac{2\pi \times \num{1.496e11}}{\num{3.156e7}}
\approx \qty{2.98e4}{m/s}$ --- trente kilomètres chaque seconde.

\textbf{17.} $\dfrac{v^2}{d} = \dfrac{(\num{2.98e4})^2}{\num{1.496e11}}
\approx \qty{5.9e-3}{N/kg}$~: le Soleil tire chaque
\omterm{def:g10:orders-of-magnitude:unit}{kilogramme} de la Terre avec environ six millinewtons.

\textbf{18.} $\dfrac{G M_{\text{Sun}}}{d^2} = \dfrac{v^2}{d}$ donne
\[
M_{\text{Sun}} = \frac{v^2 d}{G}
= \frac{\num{8.87e8} \times \num{1.496e11}}{\num{6.67e-11}}
\approx \qty{1.99e30}{kg}.
\]

\textbf{19.} $\dfrac{\num{1.99e30}}{\num{5.97e24}} \approx
\num{3.3e5}$~: le Soleil fait un tiers de million de Terres.

\textbf{20.} Parce que le même $G$ gouverne toute paire de masses, le
mesurer une fois entre deux sphères de plomb dans un hangar nous a
permis de mettre sur la balance d'abord la planète sous nos pieds,
puis l'étoile que nous \omterm{def:g10:universal-gravitation:satellite}{orbitons} --- voilà ce que
\emph{universelle} vaut.
\end{solution}
